\section{\textit{Estimation of Distribution Algorithm}}

\textit{Estimation of Distribution Algorithm} (EDA) merupakan kelas algoritma evolusioner yang menghasilkan solusi baru melalui model probabilistik, bukan operator \textit{crossover} dan mutasi seperti pada \textit{Genetic Algorithm}.
EDA, yang juga dikenal sebagai \textit{Probabilistic Model-Building Genetic Algorithm} (PMBGA)\label{first:pmbga}, mempelajari distribusi probabilitas dari individu-individu terbaik, kemudian melakukan \textit{sampling} dari distribusi tersebut untuk membentuk generasi berikutnya \parencite{larranaga_estimation_2024,lemtenneche_estimation_2023}.
Dengan demikian, pencarian solusi dipandu secara eksplisit oleh pola probabilistik yang dipelajari dari solusi berkualitas.

Secara umum, EDA menjalankan tiga tahap berulang, yaitu seleksi, estimasi, dan \textit{sampling} \parencite{larranaga_estimation_2024}.
Seleksi memilih individu terbaik berdasarkan nilai \textit{fitness}, estimasi membangun model probabilistik dari individu terpilih, sedangkan \textit{sampling} menghasilkan populasi baru dari model tersebut.
Pada persoalan diskret, distribusi dapat direpresentasikan menggunakan \textit{Bayesian network} yang memodelkan dependensi bersyarat antar-variabel secara lebih ringkas dibandingkan tabel probabilitas gabungan penuh.
Proses tersebut diulang hingga kriteria penghentian tercapai.

Berdasarkan tingkat dependensi yang dimodelkan, EDA dibedakan menjadi tiga kategori, yaitu univariat, bivariat, dan multivariat.
Eda univariat mengasumsikan seluruh variabel independen, EDA bivariat memodelkan dependensi berpasangan, sedangkan EDA multivariat memungkinkan dependensi yang lebih kompleks melalui struktur probabilistik multivariat \parencite{larranaga_estimation_2024}.

EDA telah diterapkan pada berbagai persoalan optimasi kombinatorial, termasuk masalah berbasis permutasi dengan ruang solusi besar dan kednala \textit{mutual exclusivity} antar-elemen \parencite{lemtenneche_estimation_2023}.
Meskipun ruang kebijakan ABAC bukan masalah permutasi, pembentukan kebijakan tetap bersifat kombinatorial karena setiap kandidat merupakan kombinasi diskret atribut subjek, objek, aksi, dan lingkungan.
Oleh karena itu, prinsip EDA dalam mempelajari distribusi konfigurasi yang menjanjikan berpotensi digunakan untuk mengarahkan pencarian aturan ABAC.

Namun, EDA kanonik masih menyimpan populasi berukuran $M$\label{first:label-m} secara eksplisit pada setiap generasi sehingga belum secara langsung menyelesaikan keterbatasan memori pada perangkat \textit{edge-IoT}.
Karena informasi pencarian pada EDA direpresentasikan dalam parameter model probabilistik, populasi eksplisit berpotensi digantikan oleh representasi probabilistik yang diperbarui secara inkremental.
Karakteristik ini menjadi dasar pengembangan varian EDA yang lebih \textit{compact} dengan kebutuhan memori lebih rendah.